\begin{frame}{Método Topo-Plano: Amostragem}{Análise no Domínio do Tempo}

    Multiplicamos o sinal original $x(t)$ por um trem de impulsos deslocados, assim como na amostragem ideal.

\end{frame}

\begin{frame}{Método Topo-Plano: Amostragem}{Análise no Domínio do Tempo}

    A equação para a modulação é dada por:
    \begin{equation*}
        x_p(t) = \left[ x(t) \cdot \sum_{n=-\infty}^{\infty} \delta(t - nT_s) \right] * h(t)
    \end{equation*}

    Onde $h(t)$ é o pulso retangular de duração $\tau$:
    \begin{equation*}
        h(t) =
        \begin{cases}
            1, & 0 \leq t \leq \tau \\
            0, & \text{caso contrário.}
        \end{cases}
    \end{equation*}

\end{frame}

\begin{frame}{Método Topo-Plano: Amostragem}{Interpretação Gráfica}
\centering
\begin{minipage}[c][0.66\textheight][c]{0.92\textwidth}
\centering

% ------------------------------------------------------- (1) sinal original
\only<1>{%
\begin{tikzpicture}
    \begin{axis}[pamaxis, ylabel={$x(t)$}]
        \addplot[sinalcor, line width=1pt, smooth] {\sinal};
    \end{axis}
\end{tikzpicture}\par
Sinal analógico de banda base $x(t)$.}

% ------------------------------------------------------ (2) trem de impulsos
\only<2>{%
\begin{tikzpicture}
    \begin{axis}[pamaxis, ylabel={$\delta_{T_s}(t)$}]
        \addplot[draw=none] {0};
        \pgfplotsextra{%
            \foreach \n in {0,...,\Nam}{%
                \pgfmathsetmacro\xn{\n*\Ts}%
                \draw[sinalcor, line width=1pt, -{Stealth[length=5pt,width=4pt]}]
                    (axis cs:\xn,0) -- (axis cs:\xn,0.85);
                \draw[gray, line width=.4pt] (axis cs:\xn,-0.07) -- (axis cs:\xn,0.07);}}
        \marcaTs
    \end{axis}
\end{tikzpicture}\par
Trem de impulsos com período $T_s$.}

% ------------------------------------------- (3) produto (amostragem ideal)
\only<3>{%
\begin{tikzpicture}
    \begin{axis}[pamaxis, ylabel={$x_\delta(t)$}]
        \addplot[gray!50, dashed, line width=.7pt, smooth] {\sinal};
        \pgfplotsextra{%
            \foreach \n in {0,...,\Nam}{%
                \pgfmathsetmacro\xn{\n*\Ts}%
                \pgfmathsetmacro\yn{sin(deg(2*pi*0.75*\n*\Ts + 0.4))}%
                \draw[sinalcor, line width=1pt, -{Stealth[length=5pt,width=4pt]}]
                    (axis cs:\xn,0) -- (axis cs:\xn,\yn);}}
    \end{axis}
\end{tikzpicture}\par
Produto: impulsos ponderados por $x(nT_s)$.}

% ------------------------------------------------------ (4) pulso h(t)
\only<4>{%
\begin{tikzpicture}
    \begin{axis}[pamaxis, ylabel={$h(t)$}, ymin=-0.5, ymax=1.8,
                 xlabel style={at={(axis description cs:1,0.217)}, anchor=west}]
        \addplot[draw=none] {0};
        \pgfplotsextra{%
            \draw[sinalcor, line width=1pt, fill=sinalcor!18]
                (axis cs:0,0) rectangle (axis cs:\Tau,1);
            \draw[sinalcor, line width=1pt] (axis cs:\Tau,0) -- (axis cs:1.12,0);
            \node[below=1pt, font=\small] at (axis cs:\Tau,0) {$\tau$};
            \node[left, font=\small] at (axis cs:0,1) {$1$};}
    \end{axis}
\end{tikzpicture}\par
Pulso de retenção $h(t)$, de duração $\tau$.}

% --------------------------------------------------- (5) PAM de topo plano
\only<5>{%
\begin{tikzpicture}
    \begin{axis}[pamaxis, ylabel={$x_p(t)$}]
        \addplot[gray!50, dashed, line width=.7pt, smooth] {\sinal};
        \pgfplotsextra{%
            \foreach \n in {0,...,\Nam}{%
                \pgfmathsetmacro\xa{\n*\Ts}%
                \pgfmathsetmacro\xb{\n*\Ts+\Tau}%
                \pgfmathsetmacro\yn{sin(deg(2*pi*0.75*\n*\Ts + 0.4))}%
                \draw[sinalcor, line width=.9pt, fill=sinalcor!18]
                    (axis cs:\xa,0) rectangle (axis cs:\xb,\yn);}}
        \marcaTs
    \end{axis}
\end{tikzpicture}\par
Resultado da convolução: PAM com topo plano ($d = \tau/T_s = 1/3$).}

\end{minipage}
\end{frame}

% ============================================================================
%  Efeito de abertura e deslocamento: sinal original, amostrado,
%  após o passa-baixas e após o equalizador.
%  Parâmetros e macros (prefixo \ft...) definidos no main.tex.
% ============================================================================


% ============================================================================
%  Domínio da frequência: espectro original, amostrado,
%  após o passa-baixas e após o equalizador.
% ============================================================================


% ============================================================================
%  Demonstração no domínio do tempo (3 slides)
%  Azul: etapa 1 (multiplicação pelo trem de impulsos)
%  Laranja: etapa 2 (convolução com o pulso h)
% ============================================================================
\begin{frame}{Método Topo-Plano: Amostragem}{Demonstração no Domínio do Tempo}
\small
No domínio do tempo, o processo é representado como a \ftamos{multiplicação do sinal
original por um trem de impulsos}, seguida da \ftret{convolução desse resultado com um
pulso de largura $\tau$}, conforme a equação de modulação:
\begin{equation*}
  x_p(t) = \ftamos{\Big[\, x(t)\cdot p(t) \,\Big]}\;\ftret{*\; h(t)},
  \qquad
  p(t) = \sum_{n=-\infty}^{\infty} \delta(t - nT_s),
  \qquad
  h(t) =
  \begin{cases}
    1, & 0 \leq t \leq \tau \\
    0, & \text{caso contrário.}
  \end{cases}
\end{equation*}
\begin{center}
\begin{tikzpicture}
\begin{groupplot}[group style={group name=g, group size=3 by 1, horizontal sep=1.3cm}, ftdemot]
\nextgroupplot[title={(a) Sinal original\\[1pt]{\tiny\ftstrut $x(t)$}}]
  \addplot[ftazul!40!black, thick, domain=0:\fttmax, samples=300] {cos(deg(2*pi*\ftfzero*x))};
\nextgroupplot[title={(b) Trem de impulsos\\[1pt]{\tiny\ftstrut $p(t)$}}]
  \addplot[ycomb, ftazul, thick] coordinates {\ftimpt};
  \addplot[ftimppos, ftazul] coordinates {\ftimpt};
  \draw[gray, <->, >=stealth] (axis cs:\ftTs,1.18) -- node[above, font=\tiny, inner sep=1pt]{$T_s$} (axis cs:2*\ftTs,1.18);
\nextgroupplot[title={(c) Pulso retangular\\[1pt]{\tiny\ftstrut $h(t)$}}, xmin=0, xmax=0.5,
  xtick={0,\fttau,\ftTs}, xticklabels={$0$,$\tau$,$T_s$}]
  \draw[ftlaranja!90!black, thick, fill=ftlaranja!22] (axis cs:0,0) rectangle (axis cs:\fttau,1);
\end{groupplot}
\end{tikzpicture}
\end{center}
\end{frame}

\begin{frame}{Método Topo-Plano: Amostragem}{Demonstração no Domínio do Tempo}
\small
Realizando a multiplicação, obtemos o trem de impulsos ponderado pelo sinal original,
como na amostragem ideal:
\begin{equation*}
  x_p(t) = \ftamos{\left[\, \sum_{n=-\infty}^{\infty} x(nT_s)\,\delta(t - nT_s) \right]}\;\ftret{*\; h(t)}
\end{equation*}
\begin{center}
\begin{tikzpicture}
\begin{groupplot}[group style={group name=g, group size=3 by 1, horizontal sep=1.3cm}, ftdemot]
\nextgroupplot[title={(a) Sinal original\\[1pt]{\tiny\ftstrut $x(t)$}}]
  \addplot[ftazul!40!black, thick, domain=0:\fttmax, samples=300] {cos(deg(2*pi*\ftfzero*x))};
\nextgroupplot[title={(b) Trem de impulsos\\[1pt]{\tiny\ftstrut $p(t)$}}]
  \addplot[ycomb, ftazul, thick] coordinates {\ftimpt};
  \addplot[ftimppos, ftazul] coordinates {\ftimpt};
  \draw[gray, <->, >=stealth] (axis cs:\ftTs,1.18) -- node[above, font=\tiny, inner sep=1pt]{$T_s$} (axis cs:2*\ftTs,1.18);
\nextgroupplot[title={(c) Amostragem ideal\\[1pt]{\tiny\ftstrut $x(t)\,p(t)$}}]
  \addplot[gray!55, dashed, domain=0:\fttmax, samples=300] {cos(deg(2*pi*\ftfzero*x))};
  \addplot[ycomb, ftazul, thick] coordinates {\ftamostras};
  \addplot[ftimppos, ftazul] coordinates {\ftamospos};
  \addplot[ftimpneg, ftazul] coordinates {\ftamosneg};
\end{groupplot}
\ftop{1}{2}{\times}
\ftop{2}{3}{=}
\end{tikzpicture}
\end{center}
\end{frame}

\begin{frame}{Método Topo-Plano: Amostragem}{Demonstração no Domínio do Tempo}
\small
Calculando a convolução com o pulso, obtemos o sinal amostrado final: cada impulso é
substituído por um pulso de largura $\tau$, pois $\delta(t - nT_s) * h(t) = h(t - nT_s)$:
\begin{equation*}
  \boxed{\;
  x_p(t) = \sum_{n=-\infty}^{\infty} x(nT_s)\, h(t - nT_s)
  \;}
\end{equation*}
\begin{center}
\begin{tikzpicture}
\begin{groupplot}[group style={group name=g, group size=3 by 1, horizontal sep=1.3cm}, ftdemot]
\nextgroupplot[title={(a) Amostragem ideal\\[1pt]{\tiny\ftstrut $x(t)\,p(t)$}}]
  \addplot[gray!55, dashed, domain=0:\fttmax, samples=300] {cos(deg(2*pi*\ftfzero*x))};
  \addplot[ycomb, ftazul, thick] coordinates {\ftamostras};
  \addplot[ftimppos, ftazul] coordinates {\ftamospos};
  \addplot[ftimpneg, ftazul] coordinates {\ftamosneg};
\nextgroupplot[title={(b) Pulso retangular\\[1pt]{\tiny\ftstrut $h(t)$}}, xmin=0, xmax=0.5,
  xtick={0,\fttau,\ftTs}, xticklabels={$0$,$\tau$,$T_s$}]
  \draw[ftlaranja!90!black, thick, fill=ftlaranja!22] (axis cs:0,0) rectangle (axis cs:\fttau,1);
\nextgroupplot[title={(c) Sinal amostrado\\[1pt]{\tiny\ftstrut $x_p(t)$}}]
  \addplot[gray!55, dashed, domain=0:\fttmax, samples=300] {cos(deg(2*pi*\ftfzero*x))};
  \addplot[ftlaranja!90!black, semithick, fill=ftlaranja!22] coordinates {\ftpulsos};
\end{groupplot}
\ftop{1}{2}{*}
\ftop{2}{3}{=}
\end{tikzpicture}
\end{center}
\end{frame}

% ============================================================================
%  Demonstração no domínio da frequência (3 slides)
% ============================================================================
\begin{frame}{Método Topo-Plano: Amostragem}{Demonstração no Domínio da Frequência}
\small
Para obter o espectro do sinal amostrado, aplicamos a transformada de Fourier à equação
de modulação. O produto no tempo vira convolução e a convolução vira produto:
\begin{equation*}
  x_p(t) = \ftamos{\Big[\, x(t)\cdot p(t) \,\Big]}\;\ftret{*\; h(t)}
  \quad\overset{\mathcal{F}}{\longleftrightarrow}\quad
  X_p(f) = \ftamos{\Big[\, X(f) * P(f) \,\Big]}\;\ftret{\cdot\; H(f)}
\end{equation*}
\vspace{-0.6em}
\begin{equation*}
  \ftamos{P(f) = \frac{1}{T_s}\sum_{n=-\infty}^{\infty} \delta(f - nf_s)},
  \qquad
  \ftret{H(f) = \tau\, \mathrm{sinc}(f\tau)\, e^{-j\pi f\tau}},
  \qquad
  \mathrm{sinc}(x) = \frac{\operatorname{sen}(\pi x)}{\pi x}
\end{equation*}
\begin{center}
\begin{tikzpicture}
\begin{groupplot}[group style={group name=g, group size=3 by 1, horizontal sep=1.3cm}, ftdemof]
\nextgroupplot[title={(a) Espectro original\\[1pt]{\tiny\ftstrut $X(f)$}}, yticklabels={$\frac{1}{2}$}]
  \addplot[ftimpulso, ftazul!40!black] coordinates {(-\ftfzero,1) (\ftfzero,1)};
\nextgroupplot[title={(b) Trem de impulsos\\[1pt]{\tiny\ftstrut $P(f)$}}, yticklabels={$\frac{1}{T_s}$}]
  \addplot[ftimpulso, ftazul] coordinates {\ftimpf};
\nextgroupplot[title={(c) Espectro do pulso\\[1pt]{\tiny\ftstrut $|H(f)|$}}, yticklabels={$\tau$}]
  \addplot[ftlaranja!90!black, thick, domain=-\ftfmax:\ftfmax, samples=400]
    {abs(sin(deg(pi*x*\fttau))/(pi*x*\fttau))};
  \node[font=\tiny, anchor=south west, inner sep=1pt] at (axis cs:\ftnulo,0.02) {$\frac{1}{\tau}$};
\end{groupplot}
\end{tikzpicture}
\end{center}
\end{frame}

\begin{frame}{Método Topo-Plano: Amostragem}{Demonstração no Domínio da Frequência}
\small
Ao calcular a convolução entre $X(f)$ e $P(f)$, cada impulso de $P(f)$ cria uma cópia
de $X(f)$, como na amostragem ideal:
\begin{equation*}
  \ftamos{X(f) * P(f) = \frac{1}{T_s} \sum_{n=-\infty}^{\infty} X(f - nf_s)}
\end{equation*}
\begin{center}
\begin{tikzpicture}
\begin{groupplot}[group style={group name=g, group size=3 by 1, horizontal sep=1.3cm}, ftdemof]
\nextgroupplot[title={(a) Espectro original\\[1pt]{\tiny\ftstrut $X(f)$}}, yticklabels={$\frac{1}{2}$}]
  \addplot[ftimpulso, ftazul!40!black] coordinates {(-\ftfzero,1) (\ftfzero,1)};
\nextgroupplot[title={(b) Trem de impulsos\\[1pt]{\tiny\ftstrut $P(f)$}}, yticklabels={$\frac{1}{T_s}$}]
  \addplot[ftimpulso, ftazul] coordinates {\ftimpf};
\nextgroupplot[title={(c) Amostragem ideal\\[1pt]{\tiny\ftstrut $X(f)*P(f)$}}, yticklabels={$\frac{1}{2T_s}$}]
  \addplot[ftimpulso, ftazul] coordinates {\ftrepideal};
\end{groupplot}
\ftop{1}{2}{*}
\ftop{2}{3}{=}
\end{tikzpicture}
\end{center}
\end{frame}

\begin{frame}{Método Topo-Plano: Amostragem}{Demonstração no Domínio da Frequência}
\small
Calculando, agora, a multiplicação entre o resultado da convolução e o espectro do pulso,
obtemos:
\begin{equation*}
  \boxed{\;
  X_p(f) = \frac{\tau}{T_s}\;
  \ftret{e^{-j\pi f\tau}\, \mathrm{sinc}(f\tau)}
  \;\ftamos{\sum_{n=-\infty}^{\infty} X(f - nf_s)}
  \;}
\end{equation*}
\begin{center}
\begin{tikzpicture}
\begin{groupplot}[group style={group name=g, group size=3 by 1, horizontal sep=1.3cm}, ftdemof]
\nextgroupplot[title={(a) Amostragem ideal\\[1pt]{\tiny\ftstrut $X(f)*P(f)$}}, yticklabels={$\frac{1}{2T_s}$}]
  \addplot[ftimpulso, ftazul] coordinates {\ftrepideal};
\nextgroupplot[title={(b) Espectro do pulso\\[1pt]{\tiny\ftstrut $|H(f)|$}}, yticklabels={$\tau$}]
  \addplot[ftlaranja!90!black, thick, domain=-\ftfmax:\ftfmax, samples=400]
    {abs(sin(deg(pi*x*\fttau))/(pi*x*\fttau))};
  \node[font=\tiny, anchor=south west, inner sep=1pt] at (axis cs:\ftnulo,0.02) {$\frac{1}{\tau}$};
\nextgroupplot[title={(c) Sinal amostrado\\[1pt]{\tiny\ftstrut $|X_p(f)|$}}, yticklabels={$\frac{\tau}{2T_s}$}]
  \addplot[black!55, densely dashed, domain=-\ftfmax:\ftfmax, samples=400]
    {abs(sin(deg(pi*x*\fttau))/(pi*x*\fttau))};
  \addplot[ftimpulso, ftlaranja!90!black] coordinates {\ftrepsinc};
\end{groupplot}
\ftop{1}{2}{\times}
\ftop{2}{3}{=}
\end{tikzpicture}
\end{center}
\end{frame}

% ============================================================================
%  Efeito de abertura (dinâmico: tau/T_s muda a cada clique)
% ============================================================================
\begin{frame}{Método Topo-Plano: Efeito de Abertura}{Representação e Demonstração}
\small
Após o passa-baixas ideal (ganho $T_s$, corte $f_s/2$), resta apenas a banda base,
multiplicada pelo espectro do pulso:
\begin{equation*}
  X_r'(f) = T_s\,X_p(f) = \tau\;\ftret{\mathrm{sinc}(f\tau)}\;e^{-j\pi f\tau}\;X(f),
  \qquad |f| < \frac{f_s}{2}
\end{equation*}
O fator $\ftret{\mathrm{sinc}(f\tau)}$ atenua cada componente conforme sua frequência.
Para $x(t)=\cos(2\pi\cdot 2t)$, a amplitude cai para $A = \mathrm{sinc}(2\tau)$.
\begin{center}
\foreach \d/\k in {0.25/1,0.5/2,0.75/3,1/4}{%
\only<\k>{%
\pgfmathsetmacro\tk{\d/\ftfs}%
\pgfmathsetmacro\Ak{sin(deg(pi*\ftfzero*\tk))/(pi*\ftfzero*\tk)}%
\pgfmathsetmacro\Bk{sin(deg(pi*\d/2))/(pi*\d/2)}%
\pgfmathsetmacro\dBk{20*log10(\Bk)}%
\begin{tikzpicture}
\begin{groupplot}[group style={group name=g, group size=3 by 1, horizontal sep=1.3cm}, ftdemot, height=0.18\textheight]
\nextgroupplot[title={(a) Espectro do pulso\\[1pt]{\tiny\ftstrut $|H(f)|/\tau$}},
  xmin=-\ftfmax, xmax=\ftfmax, ymin=0, ymax=1.3, xtick={-10,-5,0,5,10}, ytick={0,0.5,1},
  extra y ticks={}, xlabel={$f$ (Hz)}]
  \fill[ftroxo!8] (axis cs:-\ftfc,0) rectangle (axis cs:\ftfc,1.3);
  \addplot[ftlaranja!90!black, thick, domain=-\ftfmax:\ftfmax, samples=400]
    {abs(sin(deg(pi*x*\tk))/(pi*x*\tk))};
  \addplot[only marks, mark=*, mark size=1.5pt, ftazul] coordinates {(-\ftfzero,\Ak) (\ftfzero,\Ak)};
\nextgroupplot[title={(b) Atenuação\\[1pt]{\tiny\ftstrut em função de $\tau/T_s$}},
  xmin=0, xmax=1, ymin=0.6, ymax=1.05, xtick={0,0.25,0.5,0.75,1}, ytick={0.6,0.7,0.8,0.9,1},
  extra y ticks={}, xlabel={$\tau/T_s$}]
  \addplot[ftazul, thick, domain=0.001:1, samples=60] {sin(deg(pi*\ftfzero*x/\ftfs))/(pi*\ftfzero*x/\ftfs)};
  \addplot[ftlaranja!90!black, thick, densely dashed, domain=0.001:1, samples=60] {sin(deg(pi*x/2))/(pi*x/2)};
  \addplot[only marks, mark=*, mark size=1.6pt, ftazul] coordinates {(\d,\Ak)};
  \addplot[only marks, mark=*, mark size=1.6pt, ftlaranja!90!black] coordinates {(\d,\Bk)};
\nextgroupplot[title={(c) Amplitude no tempo\\[1pt]{\tiny\ftstrut sem o atraso, para isolar $A$}}]
  \draw[black!45, densely dashed, thin] (axis cs:0,1) -- (axis cs:\fttmax,1);
  \addplot[gray!55, dashed, domain=0:\fttmax, samples=300] {cos(deg(2*pi*\ftfzero*x))};
  \addplot[ftlaranja!90!black, thick, domain=0:\fttmax, samples=300] {\Ak*cos(deg(2*pi*\ftfzero*x))};
\end{groupplot}
\end{tikzpicture}\par\vspace{2pt}
{\scriptsize
$\tau/T_s = \pgfmathprintnumber[fixed,precision=2,use comma]{\d}$:\quad
\tikz[baseline=-0.6ex]\fill[ftazul] (0,0) circle (1.6pt);\;
$A = \mathrm{sinc}(f_0\tau) = \pgfmathprintnumber[fixed,precision=3,use comma]{\Ak}$
\qquad
\tikz[baseline=-0.6ex]\draw[ftlaranja!90!black, thick, densely dashed] (0,0) -- (0.4,0);\;
queda na borda da banda: $\mathrm{sinc}\!\left(\frac{\tau}{2T_s}\right) = \pgfmathprintnumber[fixed,precision=3,use comma]{\Bk}$
($\pgfmathprintnumber[fixed,precision=2,use comma]{\dBk}$~dB)}
}}
\end{center}
\end{frame}

% ============================================================================
%  Deslocamento (dinâmico: tau/T_s muda a cada clique)
% ============================================================================
\begin{frame}{Método Topo-Plano: Deslocamento}{Representação e Demonstração}
\small
O pulso causal é o pulso centrado atrasado de $\tau/2$. Pela propriedade do deslocamento
no tempo, $g(t-t_0) \leftrightarrow G(f)\,e^{-j2\pi f t_0}$, com $t_0 = \tau/2$:
\begin{equation*}
  h(t) = h_c\!\left(t - \frac{\tau}{2}\right)
  \;\overset{\mathcal{F}}{\longleftrightarrow}\;
  H(f) = \tau\,\mathrm{sinc}(f\tau)\;\ftret{e^{-j\pi f\tau}},
  \qquad
  \angle H(f) = -\pi f\tau
  \;\Rightarrow\;
  t_g = -\frac{1}{2\pi}\frac{d\angle H}{df} = \frac{\tau}{2}
\end{equation*}
A fase é linear: todas as frequências sofrem o mesmo atraso $\tau/2$, sem distorcer a forma do sinal.
\begin{center}
\foreach \d/\k in {0.25/1,0.5/2,0.75/3,1/4}{%
\only<\k>{%
\pgfmathsetmacro\tk{\d/\ftfs}%
\pgfmathsetmacro\atk{\d*500/\ftfs}%
\pgfmathsetmacro\pk{\ftpico+\tk/2}%
\begin{tikzpicture}
\begin{groupplot}[group style={group name=g, group size=3 by 1, horizontal sep=1.3cm}, ftdemot, height=0.18\textheight]
\nextgroupplot[title={(a) Pulso centrado e causal\\[1pt]{\tiny\ftstrut $h_c(t)$ e $h(t)$}},
  xmin=-0.15, xmax=0.28, xtick={-0.1,0,0.1,0.2}]
  \draw[gray, densely dashed, thick] (axis cs:-\tk/2,0) rectangle (axis cs:\tk/2,1);
  \draw[ftlaranja!90!black, thick, fill=ftlaranja!22] (axis cs:0,0) rectangle (axis cs:\tk,1);
  \draw[black!80, ->, >=stealth] (axis cs:0,1.12) -- (axis cs:\tk/2,1.12);
  \node[font=\tiny, anchor=south west, inner sep=1pt] at (axis cs:\tk/2,1.1) {$\tau/2$};
\nextgroupplot[title={(b) Fase do pulso\\[1pt]{\tiny\ftstrut $\angle H(f) = -\pi f\tau$}},
  xmin=-\ftfc, xmax=\ftfc, ymin=-1.8, ymax=1.8,
  xtick={-\ftfc,0,\ftfc}, xticklabels={$-\frac{f_s}{2}$,$0$,$\frac{f_s}{2}$},
  ytick={-1.5708,0,1.5708}, yticklabels={$-\frac{\pi}{2}$,$0$,$\frac{\pi}{2}$},
  extra y ticks={0}, xlabel={$f$ (Hz)}]
  \addplot[ftlaranja!90!black, thick, domain=-\ftfc:\ftfc, samples=2] {-pi*x*\tk};
\nextgroupplot[title={(c) Atraso no tempo\\[1pt]{\tiny\ftstrut sem a atenuação, para isolar $\tau/2$}}]
  \addplot[gray!55, dashed, domain=0:\fttmax, samples=300] {cos(deg(2*pi*\ftfzero*x))};
  \addplot[ftlaranja!90!black, thick, domain=0:\fttmax, samples=300] {cos(deg(2*pi*\ftfzero*(x-\tk/2)))};
  \draw[gray, densely dotted] (axis cs:\ftpico,1) -- (axis cs:\ftpico,1.18);
  \draw[ftlaranja!90!black, densely dotted] (axis cs:\pk,1) -- (axis cs:\pk,1.18);
\end{groupplot}
\end{tikzpicture}\par\vspace{2pt}
{\scriptsize
$\tau/T_s = \pgfmathprintnumber[fixed,precision=2,use comma]{\d}$:\quad
atraso $\tau/2 = \pgfmathprintnumber[fixed,precision=1,use comma]{\atk}$~ms
\qquad
\tikz[baseline=-0.6ex]\draw[gray, densely dashed, thick] (0,-0.08) rectangle (0.3,0.12);\; pulso centrado $h_c(t)$
\qquad
\tikz[baseline=-0.6ex]\draw[ftlaranja!90!black, thick, fill=ftlaranja!22] (0,-0.08) rectangle (0.3,0.12);\; pulso causal $h(t)$}
}}
\end{center}
\end{frame}

% ============================================================================
%  Reconstrução no domínio do tempo (revelado em 3 etapas)
% ============================================================================
\begin{frame}{Método Topo-Plano: Reconstrução}{Representação e Demonstração no Domínio do Tempo}
\small
Reconstruir é convoluir o sinal amostrado com $h_{PB}(t) = \mathrm{sinc}(f_s t)$, a resposta
ao impulso do passa-baixas ideal. Cada pulso vira um núcleo de interpolação
$g(t) = h(t) * h_{PB}(t)$ e, como $x(t)$ é limitado em banda, a soma é uma
\ftret{média móvel} do sinal original:
\begin{equation*}
  x_r'(t) = \sum_{n=-\infty}^{\infty} x(nT_s)\, g(t - nT_s)
  = x(t) * h(t) = \ftret{\int_{t-\tau}^{t} x(\lambda)\,d\lambda}
  \quad\xrightarrow{\text{equalizador}}\quad
  \boxed{\,x_r(t) = x\!\left(t - \frac{\tau}{2}\right)}
\end{equation*}
\begin{center}
\begin{tikzpicture}
\begin{groupplot}[group style={group name=g, group size=3 by 1, horizontal sep=1.3cm}, ftdemot, height=0.18\textheight]
\nextgroupplot[title={(a) Sinal amostrado\\[1pt]{\tiny\ftstrut $x_p(t)$}}]
  \addplot[gray!55, dashed, domain=0:\fttmax, samples=300] {cos(deg(2*pi*\ftfzero*x))};
  \addplot[ftlaranja!90!black, semithick, fill=ftlaranja!22] coordinates {\ftpulsos};
\nextgroupplot[title={(b) Núcleos de interpolação\\[1pt]{\tiny\ftstrut $x(nT_s)\,g(t-nT_s)/\tau$}}, clip=true]
  \ftnucleos
\nextgroupplot[title={(c) Soma dos núcleos\\[1pt]{\tiny\ftstrut $x_r'(t)/\tau$ e $x_r(t)$}}]
  \draw[black!45, densely dashed, thin] (axis cs:0,1) -- (axis cs:\fttmax,1);
  \addplot[gray!55, dashed, domain=0:\fttmax, samples=300] {cos(deg(2*pi*\ftfzero*x))};
  \addplot[ftroxo, thick, domain=0:\fttmax, samples=300, visible on=<2->]
    {\ftA*cos(deg(2*pi*\ftfzero*(x-\fttau/2)))};
  \addplot[ftverde, thick, densely dashed, domain=0:\fttmax, samples=300, visible on=<3->]
    {cos(deg(2*pi*\ftfzero*(x-\fttau/2)))};
\end{groupplot}
\end{tikzpicture}\par\vspace{2pt}
{\scriptsize
\tikz[baseline=-0.6ex]\draw[ftroxo, thick] (0,0) -- (0.4,0);\; após o passa-baixas
\qquad
\tikz[baseline=-0.6ex]\draw[ftverde, thick, densely dashed] (0,0) -- (0.4,0);\; após o equalizador
\qquad
\tikz[baseline=-0.6ex]\draw[gray!55, dashed] (0,0) -- (0.4,0);\; $x(t)$}
\end{center}
\pause
\end{frame}

% ============================================================================
%  Reconstrução no domínio da frequência (revelado em 3 etapas)
% ============================================================================
\begin{frame}{Método Topo-Plano: Reconstrução}{Representação e Demonstração no Domínio da Frequência}
\small
Na frequência, a reconstrução é o produto do espectro amostrado pelo passa-baixas ideal
e pelo equalizador, ambos definidos em $|f| < f_s/2$:
\begin{equation*}
  X_r(f) = X_p(f)\cdot \underbrace{T_s}_{H_{PB}(f)} \cdot\;
  \ftverde{\underbrace{\frac{1}{\tau\,\mathrm{sinc}(f\tau)}}_{H_{eq}(f)}}
  = e^{-j\pi f\tau}\,X(f)
  \quad\overset{\mathcal{F}^{-1}}{\longleftrightarrow}\quad
  \boxed{\,x_r(t) = x\!\left(t - \frac{\tau}{2}\right)}
\end{equation*}
O equalizador existe porque $\mathrm{sinc}(f\tau) \neq 0$ na banda: o primeiro nulo,
$1/\tau \geq f_s$, fica fora de $|f| < f_s/2$.
\begin{center}
\begin{tikzpicture}
\begin{groupplot}[group style={group name=g, group size=3 by 1, horizontal sep=1.3cm}, ftdemof, height=0.18\textheight]
\nextgroupplot[title={(a) Espectro amostrado\\[1pt]{\tiny\ftstrut $|X_p(f)|$ e $H_{PB}(f)$}}, yticklabels={$\frac{\tau}{2T_s}$}]
  \draw[ftroxo, densely dashed, fill=ftroxo!8] (axis cs:-\ftfc,0) rectangle (axis cs:\ftfc,1.15);
  \addplot[black!55, densely dashed, domain=-\ftfmax:\ftfmax, samples=400]
    {abs(sin(deg(pi*x*\fttau))/(pi*x*\fttau))};
  \addplot[ftimpulso, ftlaranja!90!black] coordinates {\ftrepsinc};
\nextgroupplot[title={(b) Equalizador\\[1pt]{\tiny\ftstrut $|H_{eq}(f)|$}}, yticklabels={$\frac{1}{\tau}$},
  ymax=1.8, visible on=<2->]
  \addplot[ftverde, thick, domain=-\ftfc:\ftfc, samples=100, fill=ftverde!8, visible on=<2->]
    {1/(sin(deg(pi*x*\fttau+1e-6))/(pi*x*\fttau+1e-6))} \closedcycle;
\nextgroupplot[title={(c) Sinal reconstruído\\[1pt]{\tiny\ftstrut $|X_r(f)|$}}, yticklabels={$\frac{1}{2}$}]
  \addplot[ftimpulso, ftverde, visible on=<3->] coordinates {(-\ftfzero,1) (\ftfzero,1)};
\end{groupplot}
\ftop{1}{2}{\times}
\ftop{2}{3}{=}
\end{tikzpicture}
\end{center}
\pause\pause
\end{frame}

% ============================================================================
%  Visão geral do processo completo (tempo e frequência).
%  Fica DEPOIS da reconstrução: os painéis "após o passa-baixas" e "após o
%  equalizador" só fazem sentido com as duas etapas já apresentadas.
% ============================================================================
\begin{frame}{Método Topo-Plano}{Visão Geral no Domínio do Tempo}
\centering
\begin{tikzpicture}
\begin{groupplot}[
  group style={group size=4 by 1, horizontal sep=0.85cm},
  ftaxis
]

% (a) Sinal original
\nextgroupplot[title={(a) Sinal original\\[1pt]{\tiny\ftstrut $x(t)=\cos(2\pi\cdot 2t)$}}]
  \addplot[ftazul, thick, domain=0:\fttmax, samples=400]
    {cos(deg(2*pi*\ftfzero*x))};

% (b) Sinal amostrado (topo plano)
\nextgroupplot[title={(b) Sinal amostrado\\[1pt]{\tiny\ftstrut $x_p(t)$}}]
  \addplot[gray!55, dashed, domain=0:\fttmax, samples=400]
    {cos(deg(2*pi*\ftfzero*x))};
  \addplot[ftlaranja!90!black, semithick, fill=ftlaranja!22, area legend] coordinates {\ftpulsos};
  \addplot[ycomb, ftazul, thin, mark=*, mark size=1.4pt,
    mark options={fill=ftazul, draw=white, line width=0.3pt}] coordinates {\ftamostras};

% (c) Após o passa-baixas: abertura + deslocamento
\nextgroupplot[title={(c) Após o passa-baixas\\[1pt]{\tiny\ftstrut $x_r'(t)$}},
  clip mode=individual]
  \draw[black!45, densely dashed, thin] (axis cs:0,1) -- (axis cs:\fttmax,1);
  \addplot[gray!55, dashed, domain=0:\fttmax, samples=400]
    {cos(deg(2*pi*\ftfzero*x))};
  \addplot[ftroxo, thick, domain=0:\fttmax, samples=400]
    {\ftA*cos(deg(2*pi*\ftfzero*(x-\fttau/2)))};
  \draw[gray, densely dotted] (axis cs:\ftpico,1) -- (axis cs:\ftpico,1.14);
  \draw[ftroxo, densely dotted] (axis cs:\ftpicoatr,\ftA) -- (axis cs:\ftpicoatr,1.14);
  \node[font=\tiny, anchor=south west, inner sep=1pt] at (axis cs:\ftpicoatr,1.12) {atraso $\tau/2$};

% (d) Após o equalizador: amplitude recuperada, atraso mantido
\nextgroupplot[title={(d) Após o equalizador\\[1pt]{\tiny\ftstrut $x_r(t)$}},
  clip mode=individual]
  \draw[black!45, densely dashed, thin] (axis cs:0,1) -- (axis cs:\fttmax,1);
  \addplot[gray!55, dashed, domain=0:\fttmax, samples=400]
    {cos(deg(2*pi*\ftfzero*x))};
  \addplot[ftverde, thick, domain=0:\fttmax, samples=400]
    {cos(deg(2*pi*\ftfzero*(x-\fttau/2)))};
  \draw[gray, densely dotted] (axis cs:\ftpico,1) -- (axis cs:\ftpico,1.14);
  \draw[ftverde, densely dotted] (axis cs:\ftpicoatr,1) -- (axis cs:\ftpicoatr,1.14);
  \node[font=\tiny, anchor=south west, inner sep=1pt] at (axis cs:\ftpicoatr,1.12) {atraso $\tau/2$};

\end{groupplot}
\end{tikzpicture}

\vspace{0.2em}
\mbox{\scriptsize
\tikz[baseline=-0.6ex]\draw[gray!55, dashed] (0,0) -- (0.45,0); $x(t)$\hspace{0.7em}
\tikz[baseline=-0.6ex]\draw[ftlaranja!90!black, semithick, fill=ftlaranja!22] (0,-0.08) rectangle (0.5,0.12);\hspace{3pt}pulsos de largura $\tau$\hspace{0.7em}
\tikz[baseline=-0.6ex]\fill[ftazul] (0,0) circle (1.4pt); \,amostras $x(nT_s)$\hspace{0.7em}
\tikz[baseline=-0.6ex]\draw[ftroxo, thick] (0,0) -- (0.45,0); após o passa-baixas\hspace{0.7em}
\tikz[baseline=-0.6ex]\draw[ftverde, thick] (0,0) -- (0.45,0); após o equalizador}
\end{frame}

\begin{frame}{Método Topo-Plano}{Visão Geral no Domínio da Frequência}
\centering
\begin{tikzpicture}
\begin{groupplot}[
  group style={group size=4 by 1, horizontal sep=0.85cm},
  ftfreqaxis
]

% (a) Espectro do sinal original
\nextgroupplot[title={(a) Espectro original\\[1pt]{\tiny\ftstrut $X(f)=\frac{1}{2}[\delta(f-2)+\delta(f+2)]$}}]
  \addplot[ftimpulso, ftazul] coordinates {(-\ftfzero,0.5) (\ftfzero,0.5)};

% (b) Espectro amostrado: réplicas em n f_s ponderadas pelo sinc
\nextgroupplot[title={(b) Sinal amostrado\\[1pt]{\tiny\ftstrut $X_p(f)$}}]
  \addplot[black!55, densely dashed, domain=-\ftfmax:\ftfmax, samples=500]
    {0.5*abs(sin(deg(pi*x*\fttau))/(pi*x*\fttau))};
  \addplot[ftimpulso, ftlaranja!90!black] coordinates {\ftreplicas};
  \addplot[ftimpulso, ftazul] coordinates {(-\ftfzero,\ftAmeio) (\ftfzero,\ftAmeio)};

% (c) Após o passa-baixas: só a banda base, atenuada pelo sinc
\nextgroupplot[title={(c) Após o passa-baixas\\[1pt]{\tiny\ftstrut $X_r'(f)$}}]
  \draw[black!45, densely dashed, thin] (axis cs:-\ftfmax,0.5) -- (axis cs:\ftfmax,0.5);
  \draw[ftroxo, densely dashed, fill=ftroxo!8] (axis cs:-\ftfc,0) rectangle (axis cs:\ftfc,0.6);
  \node[font=\tiny, anchor=south, inner sep=1pt, yshift=4pt, ftroxo] at (axis cs:0,0.6) {$H_{PB}(f)$};
  \addplot[ftimpulso, ftroxo] coordinates {(-\ftfzero,\ftAmeio) (\ftfzero,\ftAmeio)};

% (d) Após o equalizador: amplitude recuperada, fase linear mantida
\nextgroupplot[title={(d) Após o equalizador\\[1pt]{\tiny\ftstrut $X_r(f)$}}]
  \draw[black!45, densely dashed, thin] (axis cs:-\ftfmax,0.5) -- (axis cs:\ftfmax,0.5);
  \addplot[ftverde, densely dashed, domain=-\ftfc:\ftfc, samples=100, fill=ftverde!8]
    {0.5/(sin(deg(pi*x*\fttau+1e-6))/(pi*x*\fttau+1e-6))} \closedcycle;
  \node[font=\tiny, anchor=south, inner sep=1pt, yshift=4pt, ftverde!70!black] at (axis cs:0,0.6) {$H_{eq}(f)\propto\frac{1}{\mathrm{sinc}(f\tau)}$};
  \addplot[ftimpulso, ftverde] coordinates {(-\ftfzero,0.5) (\ftfzero,0.5)};

\end{groupplot}
\end{tikzpicture}

\vspace{0.2em}
\mbox{\scriptsize
\tikz[baseline=-0.6ex]\draw[black!55, densely dashed] (0,0) -- (0.45,0); envelope $\frac{1}{2}|\mathrm{sinc}(f\tau)|$ \hspace{0.7em}
\tikz[baseline=-0.6ex]\draw[ftazul, thick] (0,0) -- (0.45,0); banda base\hspace{0.7em}
\tikz[baseline=-0.6ex]\draw[ftlaranja!90!black, thick] (0,0) -- (0.45,0); réplicas em $nf_s \pm 2$}

\vspace{2pt}
\mbox{\scriptsize
\tikz[baseline=-0.6ex]\draw[ftroxo, densely dashed, fill=ftroxo!8] (0,-0.08) rectangle (0.5,0.12);\hspace{3pt}passa-baixas\hspace{0.7em}
\tikz[baseline=-0.6ex]\draw[ftverde, densely dashed, fill=ftverde!8] (0,-0.08) rectangle (0.5,0.12);\hspace{3pt}equalizador}
\end{frame}
